\section{III.G. Magnetic shielding of the PMTs}
The photomultipliers that are attached to the FTOF scintillators will be exposed to the combined stray magnetic fields from the CLAS12 solenoid and torus magnets. It is therefore important to study the effect of the magnetic field strength on the anode signal of a PMT and ultimately on the time resolution of the FTOF detector system. In the experimental test setup, a PMT was set on a flat, non-magnetic base in between a pair of Helmholtz coils and the anode signal of the PMT was digitized and processed by a DAQ system . This arrangement is shown in Fig. 1.

\begin{figure*}[ht]
  \includegraphics[width=2.5in]{arjun/fig_arjun_magnetic_shielding/placeholder.pdf} 
  \caption{PMT B field test setup}
  \label{fig1}
\end{figure*}

Due to the dynode-array geometry of the PMT and the electrodynamic processes which finally contribute to the signal formation, the studies need to be performed in dependence of the orientation of the PMT relative to the magnetic field. Given the axial Helmholtz magnetic field, the orientation of a PMT in it can be described by two rotational degree of freedoms; their axes are illustrated in Fig. 2 and defined by the PMT's axis of cylindrical symmetry, $\hat{z}$ and the radial axis perpendicular to $\hat{z}$ and horizontal to the flat base supporting the PMT, $\hat{\rho}$.

In the axial Helmholtz magnetic field, a PMT can be rotated around $\hat{\rho}$ until $\hat{z}$ is aligned transversely (T) or longitudinally (L) with the field direction. These are the two orientations in which a PMT is observed to have the two strongest and in a sense, independent responses; any other orientation is dominated by a superposition of the two responses. These two orientations are illustrated in Fig. 2.

\begin{figure*}[ht]
  \includegraphics[width=2.5in]{arjun/fig_arjun_magnetic_shielding/placeholder.pdf} 
  \caption{Tranverse and longitudnal orientations of PMT in axial B field}
  \label{fig2}
\end{figure*}

The additional change in the signal response of a PMT when it was rotated around $\hat{z}$ was also studied, primarily when it is already aligned in the transverse orientation (minimal to no effect due to rotations around $\hat{z}$ is observed in the longitudinal orientation) \cite{Steinman} However, after implementing the final shielding configuration (see Sec. 2.2), the PMT response no longer depends on the rotational degree of freedom around $\hat{z}$.

It is worth mentioning here that the magnetic field primarily affects the signal amplitude, while the signal shape and smoothness, which have a far larger impact on the time resolution, are mostly unaffected \cite{Steinman}. However, the loss of signal amplitude does affect the extracted time resolution, but only when the reduction reaches levels at which the Time-walk corrections become less efficient. Therefore, the signal reduction level was used as the parameter to quantify the response of a PMT in a magnetic field and its affect on the time resolution. 

Even though it was established that when the signal amplitude is reduced by 10\% in the longitudinal and the transverse orientations there is no change in the time resolution - measured to be $41\;ps$ and $38\;ps$, respectively ($39\;ps$ with no magnetic field) \cite{Steinman} -  the final design of the magnetic shielding (Sec. 2.2) is such that up to magnetic field strengths higher than those expected in CLAS12, the signal amplitude shows no reduction; the maximum stray field strength that the panel-1b PMTs are going to be exposed to is expected to be 22G at counters placed at the largest angular extent of the FTOF12 detector system, of which 2/3 (15G) will be in the axial direction \cite{CLAS12FTOFstudies} and the tests at USC were done with magnetic fields up to 30G, wholly directed in either the transverse or longitudinal directions.
